% =============================================================================
\section{La jornada final: 1 de octubre, en el piso}
% =============================================================================

\begin{frame}{El día, corrida por corrida}
  \relax{\footnotesize
  \begin{tabular}{L{1.0cm}L{3.3cm}L{9.3cm}}
    \toprule
    \textbf{Hora} & \textbf{Prueba} & \textbf{Qué pasó} \\
    \midrule
    15:21 & Aproximación 1 (oso a 90\grad) &
      25 s buscando, \textbf{cero detecciones}. Le pasó de largo al oso. \\[1pt]
    15:28 & (sin motores) &
      Con el robot quieto, 69 de 71 frames llegan en franjas. \\[1pt]
    16:02 & Aproximación de frente &
      \textbf{Primera aproximación completa}: 3,1 s. Queda a 37 cm (cinta). \\[1pt]
    16:16 & Paso de giro (20 pulsos) &
      El paso de 0,40 s renueva el 61 \% de la imagen; el de 0,25 s, el 35 \%. \\[1pt]
    16:23 & Comportamiento completo &
      Busca, encuentra, se acerca, frena. Termina torcido. \\[1pt]
    16:29 & Comportamiento completo &
      Ídem, bien centrado. Queda a $\sim$28 cm. \\[1pt]
    16:40 & Huida 1 &
      No huye: vio la mochila dos frames y el barrido giró igual. \\[1pt]
    16:43 & Huida 2 &
      No huye: la confianza de la mochila estaba al borde del umbral. \\[1pt]
    16:50 & Huida 3 &
      \textbf{Huye}, pero gira $\sim$110\grad{} y se va de costado. \\[1pt]
    16:54 & Huida 4 &
      Huida completa. Gira algo más de media vuelta. \\[1pt]
    17:07 & Demo completa 1 (90 s) &
      Oso a 2 m: tarda 36 s en verlo, lo pierde una vez, llega a los 65 s. \\[1pt]
    17:17 & Demo completa 2 (90 s) &
      Cuatro huidas ante la mochila real, tres pérdidas por mover el oso a
      mano, llegada final. \textbf{Es la del video.} \\
    \bottomrule
  \end{tabular}}

  \vspace{0.3em}
  {\footnotesize Entre corrida y corrida se cambió código: al final del día había
  60 verificaciones (eran 54), ocho parámetros con otro valor y dos nuevos.}
\end{frame}

\begin{frame}{La imagen rota (1 de 2): qué se veía}
  \begin{columns}[c,onlytextwidth]
    \begin{column}{0.32\textwidth}
      \centering
      \includegraphics[width=\textwidth]{img_sana.jpg}

      {\scriptsize Sobre el escritorio: sana. El oso con confianza 0,91 en todos
      los frames.}
    \end{column}
    \begin{column}{0.32\textwidth}
      \centering
      \includegraphics[width=\textwidth]{img_rota_37.jpg}

      {\scriptsize En el piso, \textbf{robot quieto y 12 V desenchufados}: 37 \% de
      las filas en franja. El oso está detrás.}
    \end{column}
    \begin{column}{0.32\textwidth}
      \centering
      \includegraphics[width=\textwidth]{img_rota_63.jpg}

      {\scriptsize El peor frame de esa medición: 63 \% de las filas.}
    \end{column}
  \end{columns}

  \vspace{0.5em}
  \relax{\footnotesize
  \begin{columns}[T,onlytextwidth]
    \begin{column}{0.485\textwidth}
      \begin{block}{Qué son las franjas}
        Con MJPG, cuando se pierde un paquete USB falta un tramo del JPEG, y
        el decodificador rellena con gris hasta el próximo marcador de
        reinicio. El frame llega ``entero'' pero YOLO no reconoce nada detrás.
      \end{block}
    \end{column}
    \begin{column}{0.485\textwidth}
      \begin{block}{Por qué importó tanto}
        La primera corrida no vio el oso, y la explicación a la vista era el
        paso de giro. Mirar un frame quieto antes de tocar el paso evitó
        ``arreglar'' el barrido con la cámara ciega.
      \end{block}
    \end{column}
  \end{columns}}
\end{frame}

\begin{frame}{La imagen rota (2 de 2): qué se descartó y qué no}
  \begin{columns}[T,onlytextwidth]
    \begin{column}{0.40\textwidth}
      \centering
      \includegraphics[width=0.80\textwidth]{grilla_rota.jpg}

      {\scriptsize 15:28 --- 69 de 71 frames rotos.}

      \vspace{0.3em}
      \includegraphics[width=0.80\textwidth]{grilla_limpia.jpg}

      {\scriptsize 15:52 --- 0 de 606. Nadie tocó los cables.}
    \end{column}
    \begin{column}{0.57\textwidth}
      \relax{\footnotesize
      \begin{block}{Descartado con mediciones}
        \begin{itemize}
          \item El PWM de los motores: pasa con los motores parados.
          \item La fuente de 12 V: igual con la UPS desenchufada.
          \item La carga del procesador: se rompe capturando sin YOLO.
          \item La alimentación del Pi: \cod{0x0}, 4,97 a 5,14 V.
          \item El ancho de banda USB: negocia el máximo y usa la octava parte.
          \item La cámara en sí: en los ratos buenos, 150 de 150 frames sanos.
        \end{itemize}
      \end{block}
      \begin{block}{Lo que sí se sabe}
        Son \textbf{paquetes USB perdidos}: con la traza del driver, el kernel
        registra cientos en los ratos malos y ninguno en los buenos.
      \end{block}
      \begin{alertblock}{Lo que no se sabe: la causa}
        Quedan el cable o conector de la webcam y el límite de 600 mA de los
        USB. El problema \textbf{desapareció solo} antes de poder aislarlo.
        Puede volver.
      \end{alertblock}}
    \end{column}
  \end{columns}
\end{frame}

\begin{frame}{Medir el paso de giro con la propia cámara}
  \centering
  \includegraphics[width=0.55\textwidth]{paso_giro.jpg}

  \vspace{0.2em}
  \relax{\footnotesize
  \begin{columns}[T,onlytextwidth]
    \begin{column}{0.485\textwidth}
      \begin{block}{Qué muestra}
        Un frame por pausa, con el robot quieto, a lo largo de 20 pulsos de
        giro en cuatro tandas. Comparando cada frame con el anterior
        (emparejando puntos de la mitad de arriba, que es lo lejano) sale
        cuántos píxeles se corrió la escena.
      \end{block}
    \end{column}
    \begin{column}{0.485\textwidth}
      \begin{block}{Qué se sacó}
        \begin{itemize}
          \item El paso vigente renovaba el 61 \% de la imagen: demasiado.
          \item El oso reaparece en el mismo lugar tras 15 pasos y pico: una
                vuelta son $\sim$4400 px, o sea un campo visual de
                \textbf{$\sim$52\grad}.
        \end{itemize}
      \end{block}
    \end{column}
  \end{columns}}
\end{frame}

\begin{frame}{El comportamiento completo, visto desde el robot}
  \relax{\footnotesize
  \begin{columns}[c,onlytextwidth]
    \begin{column}{0.43\textwidth}
      \centering
      \includegraphics[width=\textwidth]{aprox_a.jpg}

      \vspace{0.2em}
      \includegraphics[width=\textwidth]{aprox_b.jpg}
    \end{column}
    \begin{column}{0.55\textwidth}
      \begin{block}{Qué se ve (corrida de las 16:29)}
        Arriba a la izquierda, el último paso del barrido. El oso entra por el
        borde derecho, el robot se queda quieto (``posible objetivo''), lo
        confirma y avanza. En las filas siguientes el oso va hacia el centro
        mientras crece, hasta que el robot frena y queda detenido.
      \end{block}
      \begin{block}{Los números de esa corrida}
        15 pasos de barrido; confirmado 0,2 s después de verlo; 2,3 s de
        aproximación; el error lateral bajó de $+0{,}73$ a cero en 2 s sin
        cruzarse de lado. Cada cuadro lleva el frame, el tiempo y el motivo
        del comando.
      \end{block}
    \end{column}
  \end{columns}}
\end{frame}

\begin{frame}{La mochila que resultó ser una valija}
  \begin{columns}[T,onlytextwidth]
    \begin{column}{0.36\textwidth}
      \centering
      \includegraphics[width=\textwidth]{mochila.jpg}

      {\scriptsize Lo que ve el robot: la mochila a 1 m, desde el piso. Para
      YOLOv8n, \cod{suitcase} con confianza 0,60.}
    \end{column}
    \begin{column}{0.61\textwidth}
      \relax{\footnotesize
      \begin{block}{Cómo se supo}
        La cámara la veía perfecta y \cod{demo.sh ver} decía ``nada''.
        \cod{que_ve.py} le preguntó al modelo por \emph{todas} las clases:
        \cod{suitcase} en 15 de 15 frames, \cod{backpack} en ninguno.
      \end{block}
      \begin{block}{Por qué}
        El modelo aprendió mochilas colgadas de una espalda, vistas a la altura
        de una persona. Apoyada en el piso y mirada desde abajo, con la tapa
        enrollada, se parece más a una valija.
      \end{block}
      \begin{block}{Qué se hizo}
        \begin{itemize}
          \item \textbf{Alias}: \cod{suitcase} cuenta como la amenaza.
          \item \textbf{Umbral propio de 0,20}: a 1,2 m la confianza baja a
                0,32--0,41, justo en el umbral general de 0,35.
        \end{itemize}
      \end{block}
      \begin{alertblock}{El costo}
        Con 0,20, objetos oscuros grandes dan frames sueltos de falsa alarma:
        un marco de puerta, unos sillones negros. Uno solo no dispara nada
        (hacen falta 3 seguidos), pero hay que mirar el fondo de la sala.
      \end{alertblock}}
    \end{column}
  \end{columns}
\end{frame}

\begin{frame}{La huida, vista desde el robot}
  \relax{\footnotesize
  \begin{columns}[c,onlytextwidth]
    \begin{column}{0.43\textwidth}
      \centering
      \includegraphics[width=\textwidth]{huir_a.jpg}

      \vspace{0.2em}
      \includegraphics[width=\textwidth]{huir_b.jpg}
    \end{column}
    \begin{column}{0.55\textwidth}
      \begin{block}{Qué se ve (corrida de las 16:54)}
        Primera fila: retrocediendo, con la mochila alejándose y corriéndose
        hacia un borde (el robot gira mientras retrocede). Segunda fila: el
        giro en el lugar, con la imagen movida. De la tercera en adelante:
        avanzando, ya de espaldas, hasta quedar cerca de una pared.
      \end{block}
      \begin{block}{Los cuatro intentos}
        \begin{enumerate}
          \item No frenaba al ver la mochila sin confirmar.
          \item Confianza al borde del umbral; giraba antes de mirar.
          \item La fase 2 giraba al revés que la 1: $\sim$110\grad.
          \item Bien. Con 1,3 s se pasó de media vuelta; quedó en 1,2 s.
        \end{enumerate}
      \end{block}
    \end{column}
  \end{columns}}
\end{frame}

\begin{frame}{La demo completa (la del video)}
  \relax{\footnotesize
  \begin{columns}[c,onlytextwidth]
    \begin{column}{0.43\textwidth}
      \centering
      \includegraphics[width=\textwidth]{demo2_a.jpg}

      \vspace{0.2em}
      \includegraphics[width=\textwidth]{demo2_b.jpg}
    \end{column}
    \begin{column}{0.55\textwidth}
      \begin{block}{Qué se ve}
        Momentos de la corrida de 90 s, en un espacio abierto. Arriba: la
        mochila enfrente al arrancar (huye), y alguien moviendo el oso. Abajo:
        la mochila puesta junto al oso (caja naranja, cortada por el borde)
        dispara la huida aunque el robot vaya hacia el oso; después la mochila
        sola, y la llegada final.
      \end{block}
      \begin{block}{Cómo leer esa corrida}
        Las cuatro huidas fueron ante la mochila real. Las tres veces que
        ``perdió'' el oso, alguien lo estaba moviendo con la mano. Llegó a
        \cod{ENCONTRADO} a los 85 s, a unos 26 cm. El video está en
        \cod{media/}.
      \end{block}
    \end{column}
  \end{columns}}
\end{frame}

\begin{frame}{Lo que cambió ese día}
  \relax{\footnotesize
  \begin{tabular}{L{3.4cm}L{1.7cm}L{1.7cm}L{6.8cm}}
    \toprule
    \textbf{Qué} & \textbf{Antes} & \textbf{Después} & \textbf{Por qué} \\
    \midrule
    \cod{BUSCAR_PASO_GIRO_S} & 0,40 & 0,25 & El paso renovaba el 61 \% de la imagen \\
    \cod{BUSCAR_GIRO_S} & 9 & 17 & Una vuelta entera con pasos de 18\grad \\
    Orden del barrido & gira, mira & mira, gira & No apartar la vista antes de mirar \\
    Frenar al ver & sólo el oso & oso y mochila & La mochila no llegaba a confirmarse \\
    Envío de comandos & con el latido & en el acto & Los pulsos cortos salían con retardo al azar \\
    \cod{KP_ANG} & 0,8 & 0,5 & 0,8 se pasaba; 0,3 perdía el oso \\
    Giro al avanzar & sin límite & $\leq 0{,}7 \cdot v_{lin}$ & Ninguna rueda se detiene ni se invierte \\
    \cod{V_LIN_MAX} & 72 & 55 & Más tiempo para centrar \\
    \cod{AREA_OBJETIVO}, \cod{KP_LIN} & 0,222 y 4,5 & 0,35 y 2,4 & Frenar a 30 cm en vez de a 50 \\
    Clase de la mochila & \cod{backpack} & más \cod{suitcase} & Así la ve el modelo desde el piso \\
    Umbral de la mochila & 0,35 & 0,20 & A 1,2 m da 0,32--0,41 \\
    Fase 2 de \cod{HUIR} & a la derecha & como la fase 1 & Deshacía el giro del retroceso \\
    \cod{HUIR_GIRO_S} & 1,1 & 1,2 & 1,1 daba 110\grad; 1,3 se pasaba \\
    Búsqueda tras perder el oso & a la derecha & hacia donde lo vio & Tardó 23 s en reencontrarlo \\
    Registro & --- & columna \cod{franjas} & Saber si la cámara estaba ciega \\
    \bottomrule
  \end{tabular}}
\end{frame}
